% 20_body.tex — Day4 산출물 5/5 (⑳) 편익 실현 원장 (빈 템플릿 / 완성 예시 공용 본문)
% 드라이버: 20_편익실현원장_빈템플릿.tex (\wbblanktrue) / 20_편익실현원장_완성예시.tex (\wbblankfalse)
% 내용 출처: PM트랙/Day4_워크북.md §5.3(항목 가이드) §5.4(빈 템플릿) §5.5(온담 P1 완성 예시본)
\documentclass[10pt,a4paper]{article}
\usepackage[a4paper,margin=16mm,top=14mm,bottom=20mm]{geometry}
\def\DocNum{5}
\def\DocTitle{편익 실현 원장}
\def\DocSub{⑳ Benefits Realisation Ledger}
\def\DocVariant{\B{빈 템플릿}{온담 P1 완성 예시}}
\usepackage{wbstyle}
\begin{document}
\landscapeon
\docheader
 {\Bg{[정식 명칭과 코드]}{차세대 주문·재고 통합 플랫폼 구축 (ONE ONDAM P1)}}
 {\Bg{[v / 서명일 — 부속 관계 / 정본 관리자]}{v1.0 2027-05-21 서명 (G4 종료 보고서 부속) / 정본 재무본부(재경팀장)}}
 {\Bg{[P1 PM · 편익 소유자 실명]}{P1 PM · 오정근 · 한동석\rd{7mm}}}
 {\Bg{[재무본부 — 판정자 실명]}{정미란 CFO · 재경팀장 (판정)\rd{7mm}}}

% ------------------------------------------------------------------ 1 편익 6칸
\secbar{1. 편익 6칸 — 측정식 · 기준값 · 목표값 · 측정 시점 · 측정 책임자 실명 · 미달 시 조치 \B{}{(편익 ID·기준값·목표값·금액은 Day1 BC §4·회사 정본 §5.5; 측정식·첫 측정은 [추가 설정])}}
\guide{측정식은 \textbf{코드 수준으로}(분자·분모·산출 시스템·리포트 ID·주기). 기준값에는 반드시 \textbf{분모와 측정 시점}을 붙인다. 목표값은 값·시점. 측정 시점은 첫 측정 / G5 / M+12 / M+24. 측정 책임자 = 편익 소유자 \textbf{실명} + 측정 실무 직책(“영업본부”는 책임자가 아니다). 미달 시 조치는 원인 분류로 분기 — output이면 하자보수·유지보수, outcome이면 업무 변화 소유자, utilisation이면 채택 조치(P3). 금액 환산 편익은 \textbf{귀속}(온담 손익 vs 가맹점주)을 적는다 — 이사회 승인 금액을 온담 손익에 전액 귀속하지 않는다.}
{\scriptsize\setlength{\tabcolsep}{2pt}
\begin{longtable}{|C{12mm}|L{16mm}|L{44mm}|L{30mm}|L{22mm}|L{28mm}|L{24mm}|L{56mm}|L{22mm}|}\hline
\hd{B-ID} & \hd{편익} & \hd{측정식 (분자 · 분모 · 시스템 · 주기)} & \hd{기준값 (값 · 시점 · 분모 · 출처)} & \hd{목표값 (값 · 시점)} & \hd{측정 시점 (첫 / G5 / M+12 / M+24)} & \hd{측정 책임자 (실명 / 실무 직책)} & \hd{미달 시 조치 (원인 분류 → 누가 · 무엇을 · 언제)} & \hd{연간 금액 (승인 / 귀속)} \\ \hline \endhead
\ifwbblank
\rd{14mm}\g{B-nn} & & \g{[분자 ÷ 분모, 리포트 ID, 주기 — 외생 통제 변수]} & \g{[n\%(연도, 분모, 출처)]} & \g{[n\% (시점)]} & \g{[첫 YYYY-MM / G5 / M+12 / M+24]} & \g{[실명 / 직책]} & \g{[output → / outcome → / utilisation →]} & \g{[n / n]} \\ \hline
\rd{14mm} & & & & & & & & \\ \hline
\rd{14mm} & & & & & & & & \\ \hline
\rd{14mm} & & \g{[두 분모 병기 — 항목 2]} & \g{[A n\% (판정) / B n\% (참고)]} & \g{[A n\% / B 미설정]} & & & \g{[전제 → 이관 목록 n번]} & \\ \hline
\rd{14mm} & & & & & & & & \\ \hline
\rd{14mm} & & & & & & & & \\ \hline
\rd{14mm}\g{(참고)} & & & & & & & & \\ \hline
\rd{14mm}\g{(참고)} & & \g{[타 프로젝트 이관]} & & & & & & \\ \hline
\rd{8mm} & \textbf{직접 편익 합} & & & & & & & \g{[n / 귀속 n]} \\ \hline
\else
\textbf{B-04} & 폐기율 & 월 폐기 금액 ÷ 식자재 취급원가(C-15), 구매팀 폐기집계 + 신 플랫폼 재고 원장 폐기 사유 코드(RPT-INV-05), 월·연 누적. \textbf{수량 기준 폐기율 병기}(외생 가격 통제) & \textbf{3.9\%}(2025년 연간, 취급원가 2,530억, 구매팀 폐기집계 2025) & \textbf{2.2\%} (M+24 2028-01) & 첫 2027-06(3개월 이동 평균) / G5 2028-03 / M+12 2028-05 / M+24 2029-05 & \textbf{오정근} / 강현도 구매팀장(집계) & output(리포트 오류) → 하자보수·유지보수 / outcome(발주 정확도·수요 예측 미사용) → 오정근: 매장 발주 가이드·자동 발주 제안 채택 조치 / utilisation(매장 미준수) → P3 교육·지역매니저 점검. 미달 판정 시 오정근이 G5 전 개선 계획 제출 & 43 / \textbf{43} \\ \hline
\textbf{B-05} & 배달수수료 회피 & (배달3사 경유 매출 → 자사앱 전환 매출) × 수수료율 12.8\%(O-05), 주문 허브 채널 원장(RPT-ORD-02) + 배달 플랫폼 정산서, 월. 직영·가맹 분리 집계 & 배달 경유 242.7억(O-04), 앱 비중 2.1\%(O-03) — 2025 & 온라인 비중 18\%(B-01), M+24 & 첫 2027-09(R5 옴니 05-07 후 3개월) / G5 / M+12 / M+24 & \textbf{오정근} / 재경팀장(정산서 대조) & output(앱 연동·옴니) → 앱 벤더 UC·유지보수 / outcome(앱 전환 캠페인 미실행) → 오정근·마케팅 / 가맹 귀속분은 협의회 공유 지표로만 & 37 / \textbf{18.6}(직영 비중 50.3\%; 가맹점주 귀속 18.4) \\ \hline
\textbf{B-07} & 매장 인건비율 & 직영 매장 인건비 ÷ 직영 소매 매출(C-12), 부문별 손익 + 인사본부 급여, 분기 & \textbf{21.4\%}(2025, 2,980억) & \textbf{20.1\%} (M+24) & 첫 2027-09(2분기 마감) / G5 / M+12 / M+24 & \textbf{오정근} / 인사본부 급여 담당 & output(POS 운영 부하) → 하자보수 / outcome(인력 재배치 미실행) → 오정근·인사본부 / utilisation → P3 교육. P3 의존 명시 & 38 / \textbf{38} \\ \hline
\textbf{B-02} & 재고 정확도 & 실사 일치 SKU ÷ 실사 대상 SKU, \textbf{두 분모 병기}(StRS-034, RPT-INV-02): A 순환실사 SKU(69\%) / B 전체 SKU(100\%), 월 순환실사 + 반기 전수 & \textbf{A 91.3\%}(2025 순환실사, 판정 기준값 — 항목 2) / B 84.6\%(2026-12-18 첫 리포트, 참고) & A \textbf{98.5\%} (M+24) / B 미설정(M+12 리뷰에서 설정) & 첫 2027-06(2차 컷오버 후 센터 포함) / G5 / M+12 / M+24 & \textbf{한동석} / 물류본부 순환실사 담당(나영선 협의) & output(원장 반영 지연·KE-02·06) → 3PL 재협상·유지보수 / outcome(실사 절차 미준수) → 한동석 / 전제(검사기록) → 나영선(⑱ 이관 6) & (B-04·B-03 기여) \\ \hline
\textbf{B-03} & 결품률 & 매장 결품 SKU-일 ÷ 진열 SKU-일, 매장운영 KPI + 재고 원장 결품 플래그(RPT-INV-06), 월. \textbf{기회손실 산정식 M+6(2027-11) 확정} & \textbf{4.7\%}(2025 매장운영 KPI) & \textbf{1.5\%} (M+24) & 첫 2027-06 / 산정식 2027-11 / G5 / M+12 / M+24 & \textbf{오정근} / 영업본부 매장운영 KPI 담당 & output → 유지보수 / outcome(자동 발주 제안 미채택) → 오정근 / B-08 의존 & 미산정(M+6 확정) \\ \hline
\textbf{B-08} & 가맹 발주 리드타임 & 포털 발주 확정 시각 → 출고 확정 시각(3PL 수신 확인 기준), 가맹 388점 월 평균, 포털 로그 RPT-FRN-03 — 20:00 마감 운용 전제(CR-09) & \textbf{38시간}(2025 발주 웹 로그) & \textbf{12시간} (\textbf{M+12 2027-01 → 운영 이관 후 첫 판정 2028-01}, BC §4 시점 유지) & \textbf{첫 2027-04: 13.6시간} / 2027-10 / G5 / M+12 & \textbf{오정근} / P5 포털 운영 담당(박준영) & output(포털·배치) → 유지보수·KE-02 / outcome(3PL 출고 확정 지연 — SLA ⑥ 강등) → 한동석 3PL 재협상 / utilisation(전화·팩스 잔존 11\%) → P3 & (B-03 기여) \\ \hline
B-01 (참고) & 온라인 매출 비중 & (앱 + 배달) 매출 ÷ 소비자 대면 총매출(C-14) & 6.2\% & 18\% & B-05와 동일 & 오정근 & B-05로 환산 & (B-05) \\ \hline
B-06 (참고) & 물류 처리량 & P2 편익 계획으로 이관(Day1 BC 추가 조건 6) & 1,860 & 4,200 케이스/h & P2 G5 & 한동석 & P2 & 96 (P2) \\ \hline
& \textbf{직접 편익 합} & & & & & & & \textbf{118 / 온담 귀속 99.6} \\ \hline
\fi
\end{longtable}}

% ------------------------------------------------------------------ 2 분모 정치
\secbar[90mm]{2. 분모 정치 — 어느 기준값을 확정하는가 \B{}{(재고 정확도 기준값의 확정)}}
\guide{두 분모의 값·산출 방법·차이의 원인, 이사회 승인본의 전제(승인된 목표가 어느 분모 위에 적혔는가), 확정 기준값과 그 논리, 확장 분모의 목표 설정 시점·조건, 관련 이관 항목, 합의 서명. 유리한 분모로 바꾸지 않는다. 두 분모 중 하나를 지우지 않는다. 분모 확장이 “데이터가 없어서” 안 된다고만 적지 않고 \textbf{왜 없는지}(P1 밖의 전제)를 적어 이관 목록과 연결한다.}
{\footnotesize
\begin{tabularx}{\linewidth}{|L{30mm}|Y|}\hline
\hd{항목} & \hd{내용} \\ \hline
\B{\rd{8mm}}{}분모 A & \Bg{[범위(전체의 n\%) 기준 n\% (출처·시점)]}{순환실사 SKU(전체의 69\% — 매장 완제품·이천센터 회전 SKU) 기준 \textbf{91.3\%}(물류본부 순환실사 보고 2025, O-08)} \\ \hline
\B{\rd{8mm}}{}분모 B & \Bg{[범위(포함되는 것) 기준 n\% (출처·시점)]}{전체 SKU(100\% — 비순환 31\% 포함: 이천센터·안성 원료공장의 원료·반제품·포장재) 기준 \textbf{84.6\%}(2026-12-18 StRS-034 두 분모 리포트 첫 실행 [추가 설정])} \\ \hline
\B{\rd{10mm}}{}차이 원인 & \Bg{[왜 B의 정확도가 낮은가 — P1이 하지 않은 것과 그 이유(헌장 제외 범위·이사회 결정)]}{비순환 SKU의 입고 검수·검사기록이 종이·엑셀이라 실사 자체의 정확도가 낮고(원료 로트 단위 불일치), P1은 그 기록을 디지털화하지 않았다(헌장 §5 제외 범위 — 2025-11 이사회 반려 34억 SI 견적)} \\ \hline
\B{\rd{8mm}}{}승인본 전제 & \Bg{[이사회 승인본의 분모 명기 — 목표는 어느 분모 위의 값인가]}{이사회 승인본(회사 정본 §5.5)의 B-02 분모는 “순환실사 SKU”로 명기 — 98.5\%는 분모 A 위에 적힌 목표} \\ \hline
\B{\rd{12mm}}{}\textbf{확정 기준값} & \Bg{[A n\% → 목표 n\%(개선 n\%p)로 판정 — 논리 ① ② ③]}{\textbf{A 91.3\% → 목표 98.5\%(개선 7.2\%p)로 판정한다.} 논리: ① 승인된 목표는 승인된 분모로 판정하는 것이 정직하다 ② 분모 B로 바꾸면 개선 폭 13.9\%p(1.93배)가 되어 P1이 약속하지 않은 것을 약속한 것이 된다 ③ 분모 B의 정확도는 P1 밖의 전제(검사기록) 위에 있다} \\ \hline
\B{\rd{10mm}}{}확장 분모 목표 & \Bg{[병기·목표 미설정 — 설정 시점·조건(이관 항목·결정 회의체), 부결 시 규칙]}{B는 병기하되 목표 미설정. \textbf{M+12 리뷰(2028-05)에서 설정} — 조건: 검사기록 디지털화 과제의 이사회 재검토 결과(⑱ 이관 6번, 나영선, 2027-09 부의). 재검토 부결 시 B 목표는 “A와 동일 방법으로 측정 가능한 범위”로 한정} \\ \hline
\B{\rd{8mm}}{}합의 & \Bg{[서명자 실명·일자, 사본, 이사회 보고 문안]}{한동석·나영선·재경팀장 서명 05-21(Day3 R-12 M+3 합의의 확정), 윤태호 사본. 이사회 보고 06월에 “두 분모 병기·판정은 A” 명시} \\ \hline
\end{tabularx}}

% ------------------------------------------------------------------ 3 손익분기 · 실현 곡선
\secbar[80mm]{3. 손익분기 재확인 · 편익 실현 곡선 \B{}{(day4\_closing.py §5)}}
\guide{BC의 손익분기 실현률을 재계산(할인 관행 차이 명시)하고, \textbf{최종 AC를 넣으면 몇 \%가 되는지} 적는다. 실현 곡선(연도별 \%)과 첫 측정을 대조한다. 허용범위 $-$5\%p는 금액이 아니라 실현률 기준임을 명시하고, 판정자·회의체를 적는다. NPV만 다시 적고 손익분기를 안 적는 것이 흔한 오류다.}
{\footnotesize
\begin{tabularx}{\linewidth}{|L{46mm}|L{100mm}|Y|}\hline
\hd{항목} & \hd{값} & \hd{비고} \\ \hline
\B{\rd{9mm}}{}손익분기 편익 실현률 (BC) & \Bg{[n\% 재확인 — n억 × n = n억/년]}{\textbf{66.0\%} 재확인 — 118 × 0.66 = 77.8억/년} & \Bg{[할인율·기간·운영비 귀속 가정, NPV 범위(할인 관행)]}{할인율 8\%, 2026~2031, 운영비 27억 전액 P1 귀속(보수적). NPV(100\%) = +124(2026 할인) ~ +134(2026 무할인, BC 표기 관행)} \\ \hline
\B{\rd{9mm}}{}최종 AC 반영 & \Bg{[구축비 n → n(연도별) 시 손익분기 n\%, NPV n]}{구축비 168 → \textbf{151.95}(2026 94.0 / 2027 57.95) 시 손익분기 \textbf{62.1\%}, NPV(100\%) +138} & \Bg{[미사용 예비비의 효과 = 손익분기 n\%p 하락]}{컨틴전시 반납 4.05·관리예비비 12.0 미사용의 효과 = 손익분기 3.9\%p 하락} \\ \hline
\B{\rd{9mm}}{}실현 곡선 & \Bg{[연도별 n\% → n\%]}{2027 30\%(컷오버 후 안정화) → 2028 이후 100\%} & \Bg{[첫 측정과의 대조 — 진행률 산식]}{2027-04 B-08 13.6h(목표 12h의 진행률 (38$-$13.6)/(38$-$12) = 94\%)는 곡선보다 앞섬 — 단 B-04·B-07은 2027-09 첫 측정 전} \\ \hline
\B{\rd{10mm}}{}판정 규칙 (허용범위 편익 $-$n\%p) & \Bg{[실현률 n\% 미만 → 경고 / n\% 미만 → 중단 조건 수준 — 종료 후에는 무엇의 근거인가]}{실현률 95\% 미만 → 경고(원인 분류·개선 계획) / 66\% 미만 → 헌장 §14 중단 조건 수준(프로젝트는 종료되었으므로 프로그램 이사회의 후속 투자 판단 근거)} & \Bg{[판정자 실명·회의체]}{판정자 재무본부(정미란), 회의체 프로그램 이사회 G5} \\ \hline
\end{tabularx}}

% ------------------------------------------------------------------ 4 편익 리뷰 계획
\secbar[110mm]{4. 편익 리뷰 계획 — 첫 측정 · G5 · M+12 · M+24 \B{}{(두 시계: 프로그램 착수 2026-01 기준 M / 운영 이관 2027-05 기준)}}
\guide{리뷰 시점별 대상 항목·판정자·회의체·입력 문서·미달 시 에스컬레이션. \textbf{시계가 둘}(프로그램 착수 기준 M / 운영 이관 기준)임을 표에 명시한다 — M의 기준일이 문서마다 다르면 G5에 “언제 재는가”를 다툰다. 종료 게이트 이후의 리뷰(운영 +12·+24)를 P5·재무본부 소관으로 둔다 — 회의체가 없으면 “누가 부르는가”가 없다.}
{\footnotesize\setlength{\tabcolsep}{2.5pt}
\begin{tabularx}{\linewidth}{|L{26mm}|L{20mm}|Y|L{46mm}|L{40mm}|L{40mm}|}\hline
\hd{시점} & \hd{시계} & \hd{대상} & \hd{판정자 · 회의체} & \hd{입력} & \hd{미달 시} \\ \hline
\ifwbblank
\rd{8mm}\g{[YYYY-MM (완료)]} & \g{[운영]} & \g{[B-nn 첫 측정 값]} & \g{[실명, 공지 매체]} & & \\ \hline
\rd{8mm}\g{[YYYY-MM]} & \g{[운영 +n]} & \g{[첫 측정 항목]} & \g{[서비스 리뷰]} & \g{[리포트 ID]} & \g{[측정식 보정만]} \\ \hline
\rd{8mm} & & & & & \\ \hline
\rd{8mm} & \g{[M+n]} & \g{[산정식 확정 등 BC 추가 조건]} & & & \\ \hline
\rd{10mm}\g{\textbf{[G5 일자]}} & \g{[M+n]} & \g{[전 항목 판정 — 손익분기 n\% · 허용범위]} & \g{[재무본부 실명, 이사회]} & \g{[원장 v, n개월 실적]} & \g{[경고 → 개선 계획 / 후속 투자 판단]} \\ \hline
\rd{10mm}\g{[운영 +12]} & & \g{[원장 v2 — 확장 분모 목표 설정, 측정 조직 확인]} & \g{[재무본부·P5, 운영위]} & & \\ \hline
\rd{10mm}\g{[운영 +24]} & & \g{[원장 v3 — 유지 판정, 원장 종결 또는 상시 KPI 전환]} & & & \\ \hline
\else
2027-04 (완료) & 운영 & B-08 첫 측정 13.6h & 오정근, 협의회 공문 04-05 & 포털 로그 3월 & — \\ \hline
2027-06 & 운영 +1 & B-04·B-02·B-03 첫 측정(2차 컷오버 후 센터 포함) & 오정근·한동석, 서비스 리뷰 & RPT-INV-02·05·06 & 측정식 보정만 \\ \hline
2027-09 & 운영 +4 & B-05·B-07 첫 측정(R5 후 3개월·2분기 마감) & 오정근, 운영위 & RPT-ORD-02, 부문별 손익 & — \\ \hline
2027-11 & M+23 & B-03 기회손실 산정식 확정(BC 추가 조건 7, M+6 → 운영 +6) & 오정근·재경팀장 & 6개월 결품 데이터 & — \\ \hline
\textbf{2028-03-31 G5} & \textbf{M+26} & \textbf{전 항목 판정 — 손익분기 66\%·허용범위 $-$5\%p}(BC §4 M+24 = 2028-01 실적) & \textbf{정미란(재무본부), 프로그램 이사회(김선호)} & 원장 v1.x, 12개월 실적 & 95\% 미만 경고 → 개선 계획 / 66\% 미만 → 후속 투자 판단 \\ \hline
\textbf{2028-05} & \textbf{운영 +12} & 원장 v2 — B-02 분모 B 목표 설정, 유지 판정, 측정 조직 확인 & 재무본부·P5(박준영), 운영위 & G5 결과, 검사기록 과제 결정 & BR-04·05 대응 \\ \hline
\textbf{2029-05} & \textbf{운영 +24} & 원장 v3 — 편익 유지 판정(M+24 목표 2년 유지), 원장 종결 또는 상시 KPI 전환 & 재무본부, 운영위 & 24개월 실적 & 상시 KPI로 전환 시 원장 종결 \\ \hline
\fi
\end{tabularx}}

% ------------------------------------------------------------------ 5 ITO 책무 분리 서명표
\secbar[110mm]{5. ITO 책무 분리 서명표 — output / outcome / utilisation \B{}{(05-21, 원장 v1.0 부속)}}
\guide{편익 항목마다 세 사람 — output(플랫폼 인도: P1 PM·수행사, G4 인수로 종결), outcome(지표 달성: 편익 소유자, G5까지), utilisation(업무 변화·사용: 현업 실행 조직·P3·P5). 서명은 “책임을 진다”가 아니라 “\textbf{이 구분에 동의한다}”의 서명. 편익 책임자가 PM이거나 output과 outcome이 한 사람이면 “시스템은 됐는데 안 쓴다”가 아무의 책임도 아니게 된다. Zwikael \& Smyrk(2011).}
{\footnotesize\setlength{\tabcolsep}{2.5pt}
\begin{tabularx}{\linewidth}{|L{26mm}|L{70mm}|L{48mm}|Y|}\hline
\hd{B-ID} & \hd{output — 플랫폼 인도 (G4 인수로 종결)} & \hd{outcome — 지표 달성 (G5까지)} & \hd{utilisation — 업무 변화 · 사용} \\ \hline
\ifwbblank
\rd{9mm}\g{[B-nn 편익]} & \g{[P1 PM · 수행사 PM — 인도물 · 서명 일자]} & \g{[편익 소유자 실명 · 서명]} & \g{[실행 조직 직책·실명, P3/P5 — 서명]} \\ \hline
\rd{9mm} & & & \\ \hline
\rd{9mm} & & & \\ \hline
\rd{9mm} & & & \\ \hline
\rd{9mm} & & & \\ \hline
\rd{9mm} & & & \\ \hline
\rd{9mm}판정 & — & \g{[재무본부 실명 — 편익 환산 · G5 판정]} & \g{[원장 정본 관리자]} \\ \hline
\else
B-04 폐기율 & P1 PM·이재현(수행사) — 재고 원장·폐기 사유 코드·자동 발주 제안 인도 ✓ 05-26 & \textbf{오정근} ✓ & 매장운영팀장(발주 가이드 준수), \textbf{강현도}(폐기집계·수량 기준 병기), P3 PM(교육) ✓ \\ \hline
B-05 배달수수료 & P1 PM·이재현 — 주문 허브·앱 연동·옴니 R5 ✓ & \textbf{오정근} ✓ & 마케팅(앱 전환 캠페인), 앱 벤더 UC(박세연), 재경팀장(정산서 대조·귀속 계산) ✓ \\ \hline
B-07 인건비율 & P1 PM·이재현 — POS·매장 운영 ✓ & \textbf{오정근} ✓ & 인사본부 급여 담당(측정), 매장운영팀장(인력 재배치), P3 PM ✓ \\ \hline
B-02 재고 정확도 & P1 PM·이재현 — 통합 재고 원장·두 분모 리포트 ✓ & \textbf{한동석} ✓ & 물류본부 순환실사 담당, \textbf{나영선}(검사기록 전제·분모 B), 3PL(한동석 계약) ✓ \\ \hline
B-03 결품률 & P1 PM·이재현 — 결품 플래그·자동 발주 ✓ & \textbf{오정근} ✓ & 매장운영 KPI 담당, 재경팀장(산정식) ✓ \\ \hline
B-08 리드타임 & P1 PM·이재현 — 포털·배치 ✓ & \textbf{오정근} ✓ & \textbf{박준영}(P5 포털 운영·로그 보존), 한동석(3PL 출고 확정), P3(전화·팩스 잔존 해소) ✓ \\ \hline
판정 & — & \textbf{재무본부(정미란·재경팀장)} — 측정값의 편익 환산·G5 판정 & 원장 정본 관리 재경팀장 ✓ \\ \hline
\fi
\end{tabularx}}
\B{\small\g{output 서명의 종결 시점 / outcome·utilisation 서명의 유효 기간과 서명자 변경 시 규칙: }\hrulefill\par}{\small output 서명은 G4 인수(SAC·검수 확인)로 종결되어 P1 PM·수행사는 G5 판정에 책임이 없다. outcome·utilisation 서명은 2029-05까지 유효하며 서명자 변경 시 후임이 재서명한다(BR-04).\par}

% ------------------------------------------------------------------ 6 편익 리스크
\secbar[110mm]{6. 편익 리스크 — 귀인 · 외생변수 · 측정 조직 소멸 \B{}{(원장 부속, 자금원은 운영 예산·프로그램 예비비)}}
\guide{편익 실현을 위협하는 리스크(귀인 불가·외생변수·측정 조직 변경·분모 분쟁·전제 미충족)의 원인 → 사건 → 영향, 확률, 영향, 대응, 소유자, 검토 시점. 프로젝트 리스크 등록부 형식이되 \textbf{자금원은 컨틴전시가 아니라 운영 예산·프로그램 예비비}. 외생변수는 편익 미달의 변명이 아니라 측정식의 통제 변수로 분리한다. 프로젝트 등록부에서 이관된 편익 리스크(R-12)를 포함한다.}
{\footnotesize\setlength{\tabcolsep}{2.5pt}
\begin{tabularx}{\linewidth}{|L{30mm}|Y|C{9mm}|L{30mm}|L{62mm}|L{26mm}|L{22mm}|}\hline
\hd{BR} & \hd{원인 → 사건 → 영향} & \hd{P\%} & \hd{영향} & \hd{대응} & \hd{소유자} & \hd{검토} \\ \hline
\ifwbblank
\rd{11mm}\g{BR-01 [분모 분쟁 — R-nn 이관]} & & & & \g{[합의문·이사회 보고 문안]} & \g{[실명]} & \g{[G5]} \\ \hline
\rd{11mm}\g{BR-02 [외생변수]} & \g{[원재료가·경기 → 왜곡 → 귀인 불가]} & & & \g{[통제 변수 — 측정식에 병기]} & & \\ \hline
\rd{11mm}\g{BR-03 [귀속 분쟁]} & & & & & & \\ \hline
\rd{11mm}\g{BR-04 [측정 조직 소멸]} & \g{[담당 변경·로그 보존·서명자 이동 → 측정 단절]} & & & \g{[로그 보존 기간, 재서명 규칙, 정본 관리자]} & & \\ \hline
\rd{11mm}\g{BR-05 [전제 미충족]} & & & & & & \\ \hline
\else
BR-01 분모 분쟁 (R-12 이관) & G5에서 분모 A/B 다툼 재점화(예: A 98.5 달성·B 89) → 편익 판정 지연 & 30 & 판정 1분기 지연 & 항목 2 합의문(05-21)·이사회 보고 06월에 “판정은 A” 명시, B는 M+12 목표 설정 & 재경팀장·한동석 & G5 \\ \hline
BR-02 외생변수 (폐기율) & 원재료가 급등·이상 기후 → 취급원가 기준 폐기율 왜곡 → 귀인 불가 & 40 & B-04 43억의 ±20\% 오차 & 수량 기준 폐기율 병기(측정식), 2025 대비 원재료가 지수 통제 변수 기록 & 오정근·강현도 & 2027-06 첫 측정 \\ \hline
BR-03 B-05 귀속 분쟁 & 가맹점주 몫 18.4억을 협의회가 “본사 편익”으로 오해 → 부속서 데이터 용도 재점화(서정필) & 25 & 협의 관계 & 원장에 귀속 분리 명기, 협의회에 “가맹 귀속 지표” 공유(04-05 공문 형식) & 오정근 & 2027-09 \\ \hline
BR-04 측정 조직 소멸 & P5 포털 운영 담당 변경·포털 로그 보존 정책 부재·outcome 서명자 인사 이동 → 측정 단절 & 35 & 항목 미판정 & 로그 보존 3년(운영 문서 10종 아카이브 절차에 포함), 서명자 변경 시 재서명 규칙, 원장 정본 재무본부 & 박준영·재경팀장 & 서비스 리뷰 분기 \\ \hline
BR-05 전제 미충족 (검사기록) & 검사기록 디지털화 이사회 재검토 부결(2027-09) → 분모 B 목표 설정 불가 → 재고 정확도 편익이 순환 SKU에 한정 & 50 & B-02 확장 불가, B-04 원료 폐기 부분 미실현 & 이관 6번(나영선) 부의안에 B-02·B-04 연결 명시; 부결 시 B 목표 한정 규칙(항목 2) & 나영선·한동석 & 2027-09 이사회 \\ \hline
\fi
\end{tabularx}}
\vspace{3mm}
\noindent{\footnotesize
\begin{tabularx}{\linewidth}{|Y|Y|Y|Y|Y|}\hline
\hd{편익 소유자 — outcome (서명)} & \hd{편익 소유자 — outcome (서명)} & \hd{전제 소유자 (서명)} & \hd{판정 — 재무본부 (서명)} & \hd{output — P1 PM · 수행사 (서명)} \\ \hline
\Bg{[실명]}{오정근 2027-05-21\rd{10mm}} & \Bg{[실명]}{한동석 05-21} & \Bg{[실명]}{나영선 05-21} & \Bg{[CFO · 재경팀장]}{정미란 · 재경팀장 05-21} & \Bg{[P1 PM · 수행사 PM]}{P1 PM · 이재현 05-21} \\ \hline
\end{tabularx}}
\landscapeoff
\end{document}
